% ─────────────────────────────────────────────────────────────────────────────
%  Capítulo 5 – Representación Intermedia
% ─────────────────────────────────────────────────────────────────────────────
\chapter{Representación intermedia}
\label{chap:ir}

\section{Motivación}

La representación intermedia (\IR{}) desacopla el front-end
(específico del lenguaje fuente) del back-end (específico de la
arquitectura destino). Esto permite, en principio, añadir nuevos
lenguajes fuente o nuevas arquitecturas destino sin reescribir
el compilador completo.

\section{Formato de cuádruplas}

La \IR{} consiste en una secuencia plana de \textbf{cuádruplas} de la forma:
\[
  (op,\; \mathit{arg}_1,\; \mathit{arg}_2,\; \mathit{result})
\]

\begin{table}[H]
  \centering
  \caption{Opcodes de la representación intermedia}
  \label{tab:opcodes}
  \begin{tabular}{llp{7cm}}
    \toprule
    \textbf{Opcode}   & \textbf{Ejemplo}          & \textbf{Significado} \\
    \midrule
    \texttt{IR\_ADD}  & $(+, a, b, t_0)$         & $t_0 \leftarrow a + b$ \\
    \texttt{IR\_SUB}  & $(-, a, b, t_0)$         & $t_0 \leftarrow a - b$ \\
    \texttt{IR\_MUL}  & $(*, a, b, t_0)$         & $t_0 \leftarrow a * b$ \\
    \texttt{IR\_DIV}  & $(/, a, b, t_0)$         & $t_0 \leftarrow a / b$ \\
    \texttt{IR\_MOD}  & $(\%, a, b, t_0)$        & $t_0 \leftarrow a \bmod b$ \\
    \texttt{IR\_COPY} & $(=, a, -, x)$           & $x \leftarrow a$ \\
    \texttt{IR\_JZ}   & $(jz, cond, L, -)$       & \texttt{if} cond $= 0$ \texttt{goto} $L$ \\
    \texttt{IR\_JMP}  & $(jmp, L, -, -)$         & \texttt{goto} $L$ \\
    \texttt{IR\_LABEL}& $(label, L, -, -)$       & marca de salto $L$ \\
    \texttt{IR\_CALL} & $(call, f, -, t_0)$      & $t_0 \leftarrow f()$ \\
    \texttt{IR\_RETURN}& $(ret, a, -, -)$        & devuelve $a$ \\
    \texttt{IR\_PRINT\_INT}& $(print\_int, a, -, -)$ & imprime entero \\
    \texttt{IR\_READ\_INT} & $(read\_int, -, -, t_0)$ & lee entero de stdin \\
    \texttt{IR\_HALT} & $(halt, -, -, -)$        & fin del programa \\
    \bottomrule
  \end{tabular}
\end{table}

\section{Argumentos IR}

Cada argumento de una cuádrupla es un \texttt{IRArg} que puede ser:

\begin{itemize}
  \item \textbf{Constante entera}: \texttt{IRArg::constInt(42)}
  \item \textbf{Constante real}: \texttt{IRArg::constReal(3.14)}
  \item \textbf{Variable}: \texttt{IRArg::var("x", addr, tipo)}
  \item \textbf{Temporal}: \texttt{prog.newTemp(tipo)} — generado automáticamente
  \item \textbf{Etiqueta}: \texttt{IRArg::label("L0")}
  \item \textbf{Función}: \texttt{IRArg::func("absValue")}
  \item \textbf{Ninguno}: \texttt{IRArg::none()}
\end{itemize}

\section{Traducción de construcciones}

\subsection{Expresión aritmética}

\begin{ejemplo}
La expresión \texttt{(a + b) * c} se traduce como:
\[
  (IR\_ADD, a, b, t_0),\quad (IR\_MUL, t_0, c, t_1)
\]
\end{ejemplo}

\subsection{Sentencia if-else}

\begin{lstlisting}[style=java, caption={if-else en código fuente}]
if (x > 0) y = x; else y = -x;
\end{lstlisting}

\begin{verbatim}
(IR_GT,    x, 0,  t0)      // t0 = x > 0
(IR_JZ,    t0, L_else, -)  // if !t0 goto L_else
(IR_COPY,  x, -, y)        // then: y = x
(IR_JMP,   L_end, -, -)
(IR_LABEL, L_else, -, -)
(IR_NEG,   x, -, t1)       // t1 = -x
(IR_COPY,  t1, -, y)       // else: y = -x
(IR_LABEL, L_end, -, -)
\end{verbatim}

\subsection{Bucle while}

\begin{verbatim}
(IR_LABEL, L_start, -, -)
(IR_LTE,   i, 5, t0)
(IR_JZ,    t0, L_end, -)
... cuerpo ...
(IR_JMP,   L_start, -, -)
(IR_LABEL, L_end, -, -)
\end{verbatim}

\subsection{Llamada a función}

\begin{verbatim}
(IR_FUNC_BEGIN, absValue, -, -)
... cuerpo de absValue ...
(IR_RETURN, t3, -, -)
(IR_FUNC_END, absValue, -, -)

... en el punto de llamada ...
(IR_CALL, absValue, -, t5)  // t5 = absValue()
\end{verbatim}

\section{Visualización de la IR}

Con el flag \texttt{--ir} el compilador imprime las cuádruplas en formato
legible por el programador:

\begin{verbatim}
$ ./compiler --ir tests/backend/t04_while.java 2>&1 | head -20
[FUNC_BEGIN] sumToFive
[LABEL] L0
[LTE] i 5 -> t0
[JZ] t0 -> L1
[ADD] s i -> t1
[COPY] t1 -> s
[ADD] i 1 -> t2
[COPY] t2 -> i
[JMP] L0
[LABEL] L1
[RETURN] s
[FUNC_END] sumToFive
\end{verbatim}
